%==================================================================
% Ini adalah bab 2
% Silahkan edit sesuai kebutuhan, baik menambah atau mengurangi \section, \subsection
%==================================================================

\chapter[TINJAUAN PUSTAKA]{\\ TINJAUAN PUSTAKA}


\section{Penelitian Terdahulu}

Penelitian ini dibangun di atas fondasi penelitian yang telah ada di bidang \textit{Knowledge Graph} (KG), \textit{Natural Language Processing} (NLP), dan spesifik di bidang akademik. Tabel \ref{tab:penelitian-terdahulu} menyajikan tinjauan terhadap beberapa penelitian terdahulu yang paling relevan dan dievaluasi berdasarkan metode, temuan, serta relevansinya yang secara langsung menginformasikan desain dan implementasi dari penelitian ini.

\begin{center}
	% Definisi lebar kolom: No | Profil | Metode & Temuan | Relevansi
	\begin{longtable}{|P{0.5cm}|P{4.2cm}|P{6cm}|}
		\caption{Tinjauan Penelitian Terdahulu}\label{tab:penelitian-terdahulu}\\
		\hline
		\textbf{No} & \textbf{Profil Pustaka} & \textbf{Metode dan Temuan Utama}\\
		\hline
		\endfirsthead
		\hline
		\textbf{No} & \textbf{Profil Pustaka} & \textbf{Metode dan Temuan Utama}\\
		\hline
		\endhead
		\hline
		\multicolumn{3}{r}{\small\textit{Bersambung ke halaman berikutnya}}\\
		\endfoot
		\hline
		\endlastfoot
		
		1 &
		\textbf{Judul:} \textit{Docs2KG: Unified Knowledge Graph Construction from Heterogeneous Documents...}\newline
		\textbf{Penulis:} Qiang Sun, dkk. (2024)  &
		\textbf{Metode:} Mengembangkan kerangka kerja \textit{dual-path} (\textit{Image Converter} dan \textit{Markdown Converter}) untuk mengolah sumber data heterogen (PDF, HTML, Excel, dll.) menjadi satu KG terpadu.\newline
		\textbf{Temuan:} Efektif dalam mengintegrasikan "danau dokumen" yang beragam, memungkinkan kueri lintas format dan menyediakan fondasi untuk aplikasi \textit{knowledge-grounded RAG}.\\
		\hline
		
		2 &
		\textbf{Judul:} \textit{CourseKG: An Educational Knowledge Graph Based on Course Information for Precision Teaching}\newline
		\textbf{Penulis:} Ying Li, dkk. (2024)  &
		\textbf{Metode:} Menggunakan model \textbf{BERT-BiGRU-MHSA-CRF} untuk \textit{Named Entity Recognition} (NER) pada domain pendidikan dan ekstraksi relasi berbasis \textit{cosine similarity}.\newline
		\textbf{Temuan:} Berhasil membangun KG kurikulum yang komprehensif dari materi perkuliahan, mendukung pemetaan prasyarat dan jalur belajar personal.\\
		\hline
		
		3 &
		\textbf{Judul:} \textit{Knowledge Graph Generation and Application for Unstructured Data...}\newline
		\textbf{Penulis:} Sushmi T. Sukumar, dkk. (2024) &
		\textbf{Metode:} Membangun dan mengevaluasi \textit{pipeline} \textit{end-to-end} untuk konstruksi KG, dengan membandingkan performa model ekstraksi relasi OpenNRE vs. REBEL.\newline
		\textbf{Temuan:} Model REBEL (akurasi 87\%) terbukti jauh lebih unggul daripada OpenNRE (61.4\%). Menegaskan efektivitas pendekatan berbasis LLM generatif untuk tugas ekstraksi.\\
		\hline
		
		4 &
		\textbf{Judul:} \textit{Automatic Construction of Educational Knowledge Graphs...}\newline
		\textbf{Penulis:} Qurat Ul Ain, dkk. (2023) &
		\textbf{Metode:} Mengusulkan \textit{pipeline} tanpa supervisi untuk EduKG dengan strategi \textit{keyphrase extraction} sebelum \textit{entity linking} untuk efisiensi, serta pembobotan konsep berbasis SBERT.\newline
		\textbf{Temuan:} Strategi "ekstrak frasa kunci dulu, baru anotasi" terbukti memangkas beban komputasi. Penggunaan SBERT lebih unggul dari TF-IDF untuk pembobotan semantik.\\
		\hline
		
		5 &
		\textbf{Judul:} \textit{Building a Knowledge Graph for Recommending Experts}\newline
		\textbf{Penulis:} Behzad Rahdari \& Peter Brusilovsky (2019) &
		\textbf{Metode:} Membangun sebuah KG yang memodelkan entitas seperti Peneliti, Publikasi, dan Topik. Menggunakan KG tersebut untuk mengidentifikasi dan merekomendasikan pakar berdasarkan kueri pengguna.\newline
		\textbf{Temuan:} Sistem berbasis KG mampu memberikan rekomendasi pakar yang lebih akurat dan dapat dijelaskan dibandingkan metode tradisional. Relasi antar entitas memungkinkan penemuan pakar yang tidak terdeteksi oleh pencarian kata kunci.\\
		\hline
		
	\end{longtable}
\end{center}

\section{Landasan Teori}
\subsection{Knowledge Graph}
	\begin{subsectcontent}
		\indent \textit{Knowledge Graph} (KG) atau Graf Pengetahuan merupakan pendekatan terstruktur dalam merepresentasikan pengetahuan yang mengaitkan beragam entitas dunia nyata beserta relasinya dalam sebuah sistem yang terorganisasi. Secara matematis, KG dimodelkan sebagai graf berarah $G = (E, R, T)$, di mana $E$ adalah himpunan entitas (misalnya individu, konsep, atau peristiwa), $R$ adalah himpunan tipe relasi seperti "mengajar", "bagian\_dari", atau "berafiliasi\_dengan", dan $T \subseteq E \times R \times E$ merepresentasikan himpunan triplet fakta \cite{Hogan_2021}. Setiap triplet $(h, r, t) \in T$ merangkum satu pengetahuan: entitas kepala (\textit{head}) $h$ terhubung secara relasional dengan entitas ekor (\textit{tail}) $t$ melalui relasi $r$. Sebagai ilustrasi, triplet (Jakarta, IBUKOTA, Indonesia) atau (Prabowo, Presiden, Indonesia) mengabstraksikan relasi yang relevan dalam domain. Struktur graf pengetahuan ini menjadi dasar penalaran semantik, pencarian terarah, dan penggabungan data dari beragam sumber, sehingga KG merupakan komponen penting dalam sistem informasi dan pengambilan keputusan berbasis pengetahuan.
		
		Pembangunan KG melibatkan beberapa tahapan utama. Tahap awal adalah akuisisi pengetahuan melalui ekstraksi dari sumber teks tidak terstruktur menggunakan teknik \textit{Natural Language Processing} (NLP) seperti \textit{Named Entity Recognition} untuk mengidentifikasi berbagai entitas dan \textit{Relation Extraction} untuk menemukan relasi di antara mereka. Selain itu, tahap integrasi data terstruktur dari basis-basis pengetahuan yang sudah ada (misal Wikidata \cite{wikidata}) bertujuan memperluas cakupan graf. KG yang telah terbentuk pun bersifat dinamis dan dapat diperkaya melalui pendekatan seperti prediksi tautan (\textit{link prediction}) atau penalaran berbasis aturan untuk menginferensi fakta baru. Contohnya, aplikasi KnowEdu di bidang pendidikan memanfaatkan teknik \textit{neural sequence labeling} dan aturan asosiasi probabilistik untuk mendeteksi konsep pedagogis beserta relasi prasyarat dari data kurikulum \cite{knowedu}. Kini, dengan hadirnya Model Bahasa Besar (LLM) yang unggul dalam pemahaman semantik dan ekstraksi entitas, berbagai penelitian mutakhir memanfaatkan LLM untuk otomatisasi serta peningkatan akurasi proses pembangunan KG \cite{edge2025localglobalgraphrag}.
		
		Berbagai domain telah mengadopsi KG berkat kapabilitasnya dalam memodelkan hubungan kompleks serta integrasi data lintas sumber, misalnya:
		\begin{packed_item_sub}
			\item Search Engines dan Information Retrieval: Sistem seperti Google KG mengoptimalkan pencarian dengan memanfaatkan relasi entitas untuk menghilangkan ambiguitas kueri dan menampilkan informasi kontekstual sehingga relevansi pencarian meningkat dalam pengujian benchmark.
			\item Recommendation Systems: KG menjadi inti sistem rekomendasi modern dengan pemetaan preferensi pengguna, atribut dan keterkaitan antar item sehingga menghasilkan saran yang personal dan terjelaskan, melampaui pendekatan konvensional.
			\item Healthcare: KG mengintegrasikan data biomedis skala besar untuk mendukung proses penemuan obat, riset penyakit, serta kedokteran personal, dengan mengungkap relasi tersembunyi antara entitas biologis.
		\end{packed_item_sub}
		
		Tren penelitian terkini menunjukkan sinergi antara KG dan LLM untuk membangun sistem cerdas generasi baru. Aplikasi mutakhir mulai mengombinasikan keunggulan penyusunan relasi terstruktur dari KG dengan kemampuan penalaran dan generasi bahasa dari LLM. Misalnya, arsitektur DRAGON \cite{yasunaga2022deepbidirectionallanguageknowledgegraph} menunjukkan pencapaian state-of-the-art pada benchmark tanya jawab tingkat lanjut (QA) seperti CSQA dan OBQA dengan melatih LLM pada data KG, memungkinkan jawaban bertahap atas pertanyaan kompleks \cite{saha2018complexsequentialquestionanswering, mihaylov2018suitarmorconductelectricity}. Ke depan, arah pengembangan KG diproyeksikan menuju konstruksi graf dinamis berbasis input multimodal dan eksplorasi embedding graf untuk mendukung pengambilan keputusan secara real-time.
	\end{subsectcontent}
	
\subsection{LLM}
	\begin{subsectcontent}
		\indent Model Bahasa Besar atau biasa disebut dengan Large Language Model (LLM) merupakaan sebuah model yang bisa mengerti m
	\end{subsectcontent}

\subsection{Pemrosesan Bahasa Alami untuk Ekstraksi Pengetahuan}
	\begin{subsectcontent}
		\indent 
	\end{subsectcontent}
	
\subsection{Graf Basis Data}
\begin{subsectcontent}
	\indent 
\end{subsectcontent}

\subsection{Retrieval-Augmented Generation (RAG)}
	\begin{subsectcontent}
		\indent Retrieval-Augmented Generation (RAG) adalah pendekatan kolaboratif yang menggabungkan kemampuan model bahasa besar (LLM) yang telah dilatih sebelumnya, dengan sistem pencarian informasi eksternal untuk meningkatkan kinerja pada tugas-tugas yang membutuhkan pengetahuan mendalam. \cite{lewis2021retrievalaugmentedgenerationknowledgeintensivenlp,gao2024retrievalaugmentedgenerationlargelanguage}. Berbeda dari model bahasa tradisional yang hanya mengandalkan pengetahuan internal dalam parameter model, RAG memungkinkan sistem untuk mencari dan mengambil dokumen terkait dari database luar sebelum menghasilkan jawaban. Pendekatan ini secara efektif mampu mengatasi keterbatasan LLM terkait dengan risiko halusinasi, ppengetahuan yang usang, dan transparansi yang terbatas yang menjadikan RAG sebuah lompatan penting di dunia pemrosesan bahasa alami.
		
		\begin{figure}[!t]
			\centering
			\makebox[\textwidth][c]{\includegraphics[width=1.5\textwidth]{RAG.png}}
			\caption{Proses Retrieval-Augmented Generation untuk Question Answering \cite{gao2024retrievalaugmentedgenerationlargelanguage}}
			\label{fig:rag-process}
		\end{figure}

		Secara garis besar, framework RAG terdiri dari tiga proses utama sebagai berikut:
		\begin{enumerate}
			\item \textbf{Indexing (Indeksasi dan Ekstraksi Informasi):} Data mentah dari berbagai format (PDF, HTML, Word, Markdown) disusun dan dibersihkan agar siap  diproses.  Dokumen-dokumen ini melalui pipeline ekstraksi dan pembersihan agar seluruh informasi dapat direpresentasikan secara seragam dalam bentuk teks. Mengingat keterbatasan kontekstual LLM, teks utama akan dipotong menjadi beberapa segmen atau chunk kecil yang memudahkan analisis lanjutan. Setiap chunk kemudian dienkode ke dalam representasi vektor menggunakan model embedding tertentu dan disimpan dalam database vektor, sehingga dapat dicari kemiripannya secara efisien.
		
			\item \textbf{Retrieval (Penelusuran):} Ketika sistem menerima pertanyaan dari pengguna, query tersebut dikonversi menjadi vektor menggunakan embedding model yang sama seperti saat proses indexing. Sistem lalu mencari potongan teks dalam database yang memiliki kemiripan tinggi dengan pertanyaan. Potongan data yang paling relevan ini diambil dan digunakan sebagai dasar jawaban, memastikan jawaban yang diberikan berlandaskan sumber eksternal yang kredibel.
			\item \textbf{Generation (Generasi Jawaban):} Semua sumber relevan yang ditemukan digabung dengan pertanyaan asal, membentuk "prompt" lengkap yang dibaca oleh LLM seperti GPT untuk menghasilkan jawaban akhir. Model dapat menghasilkan jawaban dengan menggabungkan data dari dokumen yang diambil dan pengetahuan internalnya. Dalam implementasi tertentu, model bisa diatur agar hanya menggunakan dokumen yang telah diambil demi menjaga akurasi fakta. Untuk percakapan bersifat tanya-jawab berkelanjutan, histori interaksi juga bisa dimasukkan agar model konsisten menjawab.
		\end{enumerate}
		
		\FloatBarrier
		
		Dengan RAG, LLM bisa secara dinamis memanfaatkan pengetahuan eksternal membuat respons yang dihasilkan menjadi lebih akurat, informatif, dan mutakhir. Proses retrieval mengurangi risiko model "ngarang" jawaban karena jawaban merujuk pada sumber dokumen. Ini membuat RAGG sangat unggul dalam domain yang pengetahuannya berkembang cepat, seperti berita, publikasi sains, ataupun regulasi yang berubah secara periodik. Pengguna juga bisa menelusuri dari mana sumber jawaban AI berasal, meningkatkan transparansi dan kepercayaan
		
		Sebagai ilustrasi praktis, misal pengguna ingin mengetahui perkembangan suatu topik berita yang sedang viral beberapa hari terakhir. Model LLM berbasis pre-trained umumnya kurang mampu memberikan informasi terbaru, sehingga respons yang dihasilkan sering kurang relevan dengan konteks aktual. Namun, dengan penerapan skema RAG seperti pada \figref{fig:rag-process}, sistem akan secara otomatis mengambil dan mengintegrasikan dokumen berita terkini sebagai referensi utama, sehingga jawaban yang dihasilkan lebih faktual dan sesuai kebutuhan pengguna.

		Seiring dengan meningkatnya kebutuhan akan informasi yang akurat dan terkini, RAG telah diintegrasikan dan diadopsi di berbagai bidang aplikasi dalam ekosistem teknologi modern:
		\begin{subsectlist}
			\begin{packed_item}
				\item \textbf{Open-Domain Question Answering:} RAG terbukti meningkatkan performa QA tanya jawab terbuka karena mampu mengambil referensi dari Wikipedia atau dataset yang relevan, dan secara konsisten melampaui model ekstraktif berbasis machine reading comprehension pada berbagai benchmark seperti pada Natural Questions dan TriviaQA.
				\item \textbf{Aplikasi Legal dan Medis:} Di domain hukum serta kesehatan, RAG digunakan untuk penelusuran dokumen kasus hukum terbaru atau artikel ilmiah medis guna memberi landasan keputusan berbasis data mutakhir dan referensi valid.
				\item \textbf{Chatbot dan Virtual Assistant:} Integrasi RAG pada sistem percakapan seperti ChatGPT memungkinkan respons lebih akurat karena mampu mengambil data keuangan, dokumentasi teknis, ataupun berita terkini secara real-time.
			\end{packed_item}
		\end{subsectlist}
	\end{subsectcontent}

\subsection{Graph based RAG}
\begin{subsectcontent}
	\indent \textit{Graph-RAG} adalah sebuah evolusi dari arsitektur RAG yang dirancang untuk meningkatkan mekanisme pengambilan informasi dengan memanfaatkan sumber pengetahuan terstruktur, seperti \textit{Knowledge Graph} (KG) atau basis data relasional \cite{peng2024graphretrievalaugmentedgenerationsurvey}. Berbeda dengan RAG konvensional yang umumnya mengambil potongan teks tidak terstruktur, Graph-RAG secara spesifik mengambil elemen-elemen graf yang relevan—seperti entitas, relasi, atau bahkan seluruh subgraf—untuk memperkaya konteks yang diberikan kepada Model Bahasa Besar (LLM).
	
	Pendekatan ini menghasilkan beberapa keunggulan signifikan seperti, penalaran yang lebih presisi, akurasi faktual yang lebih tinggi, serta \textbf{interpretasi dan transparansi} yang lebih baik, karena respons yang dihasilkan dapat dilacak kembali ke entitas dan relasi spesifik di dalam graf. Hal ini menjadikan Graph-RAG sangat ideal untuk aplikasi yang menuntut konsistensi logika dan pemahaman mendalam terhadap hubungan antar-entitas, seperti pada sistem rekomendasi akademik yang akan dikembangkan dalam penelitian ini.
\end{subsectcontent}

\subsection{Konsep Dasar Rangkaian Digital}
Rangkaian digital beroperasi dengan sinyal diskrit, cocok untuk pemrosesan informasi digital. Bagian ini menguraikan prinsip dasar rangkaian digital serta komponen-komponen utama yang mendukung fungsi-fungsi digital dalam proyek ini.

\subsection{Teknik Pengolahan Sinyal pada Sistem Elektronika}
Pengolahan sinyal adalah proses penting untuk interpretasi informasi dari lingkungan. Bagian ini membahas metode umum dalam pengolahan sinyal yang diterapkan pada sistem elektronika, termasuk teknik yang digunakan dalam pemfilteran atau pemrosesan data.

\subsection{Pengkabelan dan Pengaturan Sirkuit untuk Keandalan Sistem}
Pengkabelan dan tata letak yang baik meningkatkan keandalan sistem secara keseluruhan. Bagian ini menguraikan teknik pengkabelan dan pengaturan sirkuit yang efektif, serta bagaimana hal ini dapat mempengaruhi performa sistem.

\subsection{Pengendalian dan Penggerak (Motor Driver, Relay, dsb.)}
Bagian ini menjelaskan penggunaan penggerak seperti motor driver dan relay untuk menggerakkan komponen mekanis. Diperlukan teknik pengendalian khusus untuk memastikan bahwa setiap penggerak bekerja sesuai dengan tujuan sistem.

\section{Teknologi yang Digunakan}
Bagian ini mengulas teknologi yang umum digunakan dalam proyek berbasis elektronika, seperti mikrokontroler, sensor, aktor, dan teknologi komunikasi nirkabel. Teknologi ini memungkinkan sistem untuk merespons lingkungan dan berinteraksi dengan pengguna.

\subsection{Mikrokontroler dan Mikroprosesor}
Mikrokontroler dan mikroprosesor berfungsi sebagai unit pemrosesan utama dalam sistem elektronika. Bagian ini menguraikan arsitektur dasar, bahasa pemrograman yang relevan, serta protokol komunikasi yang sering digunakan dalam proyek ini.

\subsection{Sensor dan Aktuator}
Sensor dan aktuator memungkinkan interaksi sistem dengan lingkungannya. Bagian ini membahas jenis-jenis sensor yang digunakan, cara kerja, dan integrasinya ke dalam sistem agar sistem dapat mengumpulkan data dan merespons secara aktif.

\subsection{Teknologi Nirkabel}
Teknologi nirkabel seperti Bluetooth dan Wi-Fi memungkinkan komunikasi jarak jauh dalam sistem IoT. Bagian ini menguraikan jenis teknologi nirkabel yang relevan, termasuk protokol komunikasi dan aspek keamanan yang perlu dipertimbangkan.

\section{Metode Kontrol dan Kecerdasan Buatan}
Bagian ini membahas metode kontrol dan kecerdasan buatan yang diterapkan dalam sistem elektronika untuk mencapai pengendalian yang lebih cerdas dan otomatis, seperti kontrol PID, logika fuzzy, dan deep learning.

\subsection{Pengendalian PID (Proportional-Integral-Derivative)}
PID adalah metode kontrol yang efektif dalam mengatur respons sistem. Bagian ini menjelaskan prinsip dasar PID, aplikasinya dalam pengaturan sistem elektronika, serta teknik tuning yang dapat meningkatkan stabilitas dan respons sistem.

\subsection{Fuzzy Logic Control}
Logika fuzzy adalah metode kontrol fleksibel yang sering digunakan dalam sistem nonlinear. Bagian ini menjelaskan konsep dasar logika fuzzy, serta cara implementasi dan manfaatnya dalam pengendalian sistem yang kompleks.

\subsection{Deep Learning}
Deep learning memungkinkan sistem untuk belajar dari data, yang sangat berguna dalam aplikasi otomatisasi. Bagian ini menguraikan algoritma dasar dalam deep learning, seperti CNN dan RNN, serta penerapannya dalam pengembangan sistem IoT.

\subsection{Perbandingan dan Pemilihan Metode yang Sesuai}
Bagian ini membahas perbandingan antara berbagai metode kontrol yang tersedia, menjelaskan kelebihan dan kekurangannya masing-masing, serta bagaimana memilih metode yang paling sesuai untuk aplikasi proyek ini.

\section{Konsep Engineering Design Process}
Engineering Design Process adalah metodologi sistematis yang digunakan untuk merancang sistem secara efektif. Bagian ini menguraikan prinsip utama dan langkah-langkah dari Engineering Design Process dalam konteks pengembangan sistem elektronika.

\subsection{Pengertian dan Langkah-langkah Engineering Design Process}
Engineering Design Process adalah proses desain iteratif yang mencakup beberapa tahapan untuk mencapai desain yang optimal. Bagian ini menjelaskan langkah-langkah dasar yang terlibat dan bagaimana proses ini diadaptasi dalam proyek ini.

\subsection{Aplikasi Engineering Design Process pada Proyek Elektronika}
Bagian ini membahas penerapan Engineering Design Process dalam proyek elektronika untuk mencapai hasil desain yang optimal. Diuraikan langkah-langkah praktis dalam menerapkan metodologi ini.

\subsection{Studi Kasus Implementasi Engineering Design Process dalam Desain Elektronika}
Studi kasus ini menunjukkan contoh penerapan Engineering Design Process dalam desain sistem yang relevan dengan proyek. Dengan studi kasus ini, pembaca dapat memahami implementasi proses desain secara nyata.

\subsection{Teknik Evaluasi dan Optimasi Desain}
Teknik evaluasi dan optimasi desain merupakan langkah penting dalam proses desain yang berkelanjutan. Bagian ini menjelaskan metode yang digunakan untuk mengevaluasi dan menyempurnakan desain agar mencapai hasil terbaik.

\section{Penelitian Terdahulu yang Relevan}
Bagian ini berisi ulasan terhadap penelitian terdahulu yang relevan dengan proyek ini. Tujuannya adalah untuk melihat pendekatan yang telah digunakan, menemukan kelebihan dan kekurangannya, serta mengidentifikasi inovasi yang dapat dikembangkan.

\subsection{Tinjauan Penelitian Terdahulu tentang Proyek Serupa}
Bagian ini membahas penelitian terdahulu yang serupa dengan proyek ini. Tinjauan ini bertujuan untuk memahami bagaimana proyek ini dapat memberikan kontribusi yang berbeda atau lebih baik.

\subsection{Analisis Kekurangan dan Kelebihan Metode pada Penelitian Terdahulu}
Setiap metode yang digunakan dalam penelitian terdahulu memiliki kelebihan dan kekurangan. Bagian ini mengidentifikasi aspek yang perlu diperbaiki atau dikembangkan lebih lanjut berdasarkan analisis metode-metode tersebut.

\subsection{Inovasi dan Kontribusi yang Dibawa dalam Penelitian Ini}
Penelitian ini membawa inovasi tertentu yang berkontribusi dalam memperkaya hasil penelitian sebelumnya. Bagian ini menjelaskan kontribusi utama proyek ini terhadap bidang elektronika, serta perbedaan yang ditawarkan.